\documentclass[../../../include/open-logic-section]{subfiles}

\begin{document}

\olfileid[hi]{sth}{ordinals}{basic}
\olsection{क्रमसूचकों के मूल गुण}

हमने देखा कि पहले कुछ क्रमसूचक प्राकृतिक संख्याएँ हैं। क्रमसूचक सिद्धांत
विकसित करने का मुख्य कारण प्राकृतिक संख्याओं पर लागू आगमन के सिद्धांत का
विस्तार करना है। प्राथमिक परिणामों की एक शृंखला से हम वहाँ तक पहुँचेंगे।

\begin{lem}\ollabel{ordmemberord}
क्रमसूचक का प्रत्येक !!{element} क्रमसूचक है।
\end{lem}

\begin{proof}
$\alpha$ को ऐसा क्रमसूचक मानें कि $b \in \alpha$। चूँकि $\alpha$
संक्रमणशील है, $b \subseteq \alpha$। अतः जैसे $\in$, $\alpha$ को
सुव्यवस्थित करता है, वैसे ही वह $b$ को भी सुव्यवस्थित करता है।

$b$ की संक्रमणशीलता देखने के लिए मान लें $x \in c \in b$। चूँकि
$b \subseteq \alpha$, इसलिए $c \in \alpha$। फिर, चूँकि $\alpha$
संक्रमणशील है, $c \subseteq \alpha$, जिससे $x \in \alpha$। अतः
$x, c, b \in \alpha$। किंतु $\in$, $\alpha$ को सुव्यवस्थित करता है,
इसलिए \olref[sth][ordinals][wo]{wo:strictorder} से $\in$, $\alpha$ पर
संक्रमणशील संबंध है। अतः $x \in c \in b$ से $x \in b$। सामान्यीकरण करने
पर $c \subseteq b$।
\end{proof}

\begin{cor}\ollabel{ordissetofsmallerord}
किसी भी क्रमसूचक~$\alpha$ के लिए
$\alpha = \Setabs{\beta \in \alpha}{\beta \text{ क्रमसूचक है}}$।
\end{cor}

\begin{proof}
\olref{ordmemberord} से तत्काल।
\end{proof}

अगले दो मुख्य परिणामों, \olref{ordinductionschema} और
\olref{ordtrichotomy}, का मोटा आशय यह है कि स्वयं क्रमसूचक सदस्यता द्वारा
सुव्यवस्थित हैं:

\begin{thm}[परिमितेतर आगमन]\ollabel{ordinductionschema}
किसी भी सूत्र $\phi(x)$ के लिए:
\[
	\text{यदि }\exists \alpha \phi(\alpha)\text{, तो }\exists \alpha(\phi(\alpha)
	\land  (\forall \beta \in \alpha) \lnot \phi(\beta))
\]
जहाँ प्रदर्शित परिमाणक निहित रूप से क्रमसूचकों तक सीमित हैं।
\end{thm}

\begin{proof}
%\olref{ordinductionschema1}. 
किसी क्रमसूचक $\alpha$ के लिए $\phi(\alpha)$ मान लें। यदि
$(\forall \beta \in \alpha) \lnot \phi(\beta)$, तो काम पूरा है। अन्यथा,
चूँकि $\alpha$ क्रमसूचक है, उसका कोई $\in$-लघुतम !!{element} है जो $\phi$
को संतुष्ट करता है; और \olref{ordmemberord} से वह क्रमसूचक है।
%	\olref{ordinductionschema2}. Suppose there is some ordinal
%	$\gamma$ such that $\lnot\phi(\gamma)$. Then by
%	\olref{ordinductionschema1}, there is an $\in$-minimal ordinal
%	$\alpha$ for which $\lnot\phi(\alpha)$. So $(\forall \beta <
%	\alpha) \phi(\beta)$, rendering the antecedent of the conditional
%	false.
\end{proof}
\noindent
ध्यान दें कि हम \olref{ordinductionschema} को समान रूप से इस योजना में
व्यक्त कर सकते हैं:
\[
\text{यदि }\forall \alpha((\forall \beta \in \alpha)\phi(\beta) \lif
\phi(\alpha))\text{, तो }\forall \alpha\phi(\alpha)
\]
इसके लिए \olref{ordinductionschema} में बस $\lnot\phi(\alpha)$ लेकर प्राथमिक
तार्किक रूपांतरण करने होते हैं।

\begin{thm}[त्रिकोटी]\ollabel{ordtrichotomy}
किन्हीं क्रमसूचकों $\alpha$ और $\beta$ के लिए
$\alpha \in \beta \lor \alpha = \beta \lor \beta \in \alpha$।
\end{thm}

\begin{proof}
प्रमाण दोहरे आगमन द्वारा है, अर्थात \olref{ordinductionschema} का दो बार
उपयोग करके। कहें कि $x$, $y$ से \emph{तुलनीय} है तभी और केवल तभी जब
$x \in y \lor x = y \lor y \in x$।

आगमन के लिए मान लें कि $\alpha$ में प्रत्येक क्रमसूचक \emph{प्रत्येक}
क्रमसूचक से तुलनीय है। आगे के आगमन के लिए मान लें कि $\alpha$, $\beta$ में
प्रत्येक क्रमसूचक से तुलनीय है। हम दिखाएँगे कि $\alpha$, $\beta$ से तुलनीय
है। $\beta$ पर आगमन से निष्पन्न होगा कि $\alpha$ प्रत्येक क्रमसूचक से तुलनीय
है; और फिर $\alpha$ पर आगमन से \emph{प्रत्येक} क्रमसूचक \emph{प्रत्येक}
क्रमसूचक से तुलनीय है, जैसा अपेक्षित था। $\alpha \notin \beta$ और
$\beta \notin \alpha$ मानकर $\alpha = \beta$ दिखाना पर्याप्त है।

$\alpha \subseteq \beta$ दिखाने के लिए $\gamma \in \alpha$ को नियत करें;
\olref{ordmemberord} से यह क्रमसूचक है। अतः पहली आगमन परिकल्पना से $\gamma$,
$\beta$ से तुलनीय है। किंतु यदि या तो $\gamma = \beta$ या
$\beta \in \gamma$, तो (आवश्यक होने पर $\alpha$ की संक्रमणशीलता का उपयोग
करके) $\beta \in \alpha$, जो हमारी मान्यता के विरुद्ध है; इसलिए
$\gamma \in \beta$। सामान्यीकरण करने पर $\alpha \subseteq \beta$।

दूसरी आगमन परिकल्पना का उपयोग करने वाला ठीक समान तर्क
$\beta \subseteq \alpha$ दिखाता है। अतः $\alpha = \beta$।
\end{proof}\noindent
इसलिए हम कभी-कभी $\alpha \in \beta$ के बजाय $\alpha <\beta$ लिखेंगे,
क्योंकि $\in$ क्रम संबंध जैसा व्यवहार कर रहा है। परिचय और यह तथ्य कि
$\alpha \in \beta \lor \alpha = \beta$ की अपेक्षा $\alpha \leq \beta$
लिखना सरल है, इसके अतिरिक्त कोई गहरा कारण नहीं।\footnote{हम
$\alpha \mathrel{\underline{\in}} \beta$ लिख सकते थे; किंतु वह पूरी तरह
अमानक होता।}

हमारे पिछले परिणामों के दो त्वरित परिणाम ये हैं, जिनमें पहला हमारे नए
संकेतन को काम में लाता है:

\begin{cor}\ollabel{ordordered}
यदि $\exists \alpha\phi(\alpha)$, तो $\exists \alpha(\phi(\alpha)
\land \forall \beta(\phi(\beta) \lif \alpha \leq \beta))$.
इसके अतिरिक्त किन्हीं क्रमसूचकों $\alpha, \beta, \gamma$ के लिए
$\alpha \notin \alpha$ और
$\alpha \in \beta \in \gamma \lif \alpha \in \gamma$, दोनों।
\end{cor}

\begin{proof}
ठीक \olref[wo]{wo:strictorder} की तरह।
\end{proof}

\begin{prob}
\olref[sth][ordinals][basic]{ordtrichotomy} के प्रमाण का ``ठीक समान तर्क''
पूर्ण कीजिए।
\end{prob}

\begin{cor}\ollabel{corordtransitiveord}
$A$ क्रमसूचक तभी और केवल तभी है जब $A$ क्रमसूचकों का संक्रमणशील समुच्चय है।
\end{cor}

\begin{proof}
\emph{बाएँ से दाएँ।} \olref{ordmemberord} से। \emph{दाएँ से बाएँ।}
यदि $A$ क्रमसूचकों का संक्रमणशील समुच्चय है, तो
\olref{ordinductionschema} और \olref{ordtrichotomy} से $\in$, $A$ को
सुव्यवस्थित करता है।
\end{proof}

हमने \olref{ordinductionschema} और \olref{ordtrichotomy} का आशय यह बताया
कि $\in$ क्रमसूचकों को सुव्यवस्थित करता है। किंतु निम्न परिणाम के कारण ऐसे
दावे के विषय में हमें \emph{बहुत सावधान} रहना होगा:

\begin{thm}[बुराली--फ़ोर्टी विरोधाभास]\ollabel{buraliforti}
समस्त क्रमसूचकों का कोई समुच्चय नहीं है।
\end{thm}

\begin{proof}
विरोधाभास द्वारा प्रमाण के लिए मान लें कि $O$ समस्त क्रमसूचकों का समुच्चय
है। यदि $\alpha \in \beta \in O$, तो \olref{ordmemberord} से $\alpha$
क्रमसूचक है, अतः $\alpha \in O$। इसलिए $O$ संक्रमणशील है और इस प्रकार
\olref{corordtransitiveord} से $O$ क्रमसूचक है। फलतः $O \in O$, जो
\olref{ordordered} के विरुद्ध है।
\end{proof}
\noindent
इस परिणाम का नाम \citeauthor{Burali-Forti1897} पर है। किंतु 1899 में
डेडेकिंड को लिखे पत्र में कैंटर ने पहली बार स्पष्टतः वह \emph{विरोधाभास}
देखा जो समस्त क्रमसूचकों का समुच्चय मानने से उत्पन्न होता है। वैन हेयेनॉर्ट
के शब्दों में:
\begin{quote}
  स्वयं बुराली--फ़ोर्टी ने माना कि विरोधाभास \emph{असंगति तक अपचय} द्वारा
  यह परिणाम स्थापित करता है कि क्रमसूचकों का स्वाभाविक क्रम केवल एक आंशिक
  क्रम है।
  \citep[p.~105]{Heijenoort1967}
\end{quote}
बुराली--फ़ोर्टी की भूल को अलग रखकर पिछले विवेचन का सार इस प्रकार है।
क्रमसूचक ऐसे समुच्चय हैं जो अलग-अलग सदस्यता द्वारा सुव्यवस्थित हैं और सामूहिक
रूप से भी सदस्यता द्वारा सुव्यवस्थित हैं (बिना सामूहिक रूप से कोई समुच्चय
बनाए)।

इसे पूर्ण करने के लिए क्रमसूचकों के कुछ और मूल गुण ये हैं, जो
\olref{ordinductionschema} और \olref{ordtrichotomy} से निष्पन्न होते हैं।

\begin{prop}
क्रमसूचकों का प्रत्येक कड़ाई से अवरोही अनुक्रम परिमित है।
\end{prop}

\begin{proof}
क्रमसूचकों के किसी अनंत कड़ाई से अवरोही अनुक्रम
$\alpha_0 > \alpha_1 > \alpha_2 > \ldots$ का कोई $<$-न्यूनतम सदस्य नहीं
है, जो \olref{ordinductionschema} के विरुद्ध है।
\end{proof}

\begin{prop}\ollabel{ordinalsaresubsets}
किन्हीं क्रमसूचकों $\alpha, \beta$ के लिए
$\alpha \subseteq \beta \lor \beta \subseteq \alpha$।
\end{prop}

\begin{proof}
यदि $\alpha \in \beta$, तो $\beta$ की संक्रमणशीलता से
$\alpha \subseteq \beta$। इसी प्रकार यदि $\beta \in \alpha$, तो
$\beta \subseteq \alpha$। और यदि $\alpha = \beta$, तो
$\alpha \subseteq \beta$ और $\beta \subseteq \alpha$। अतः
\olref{ordtrichotomy} से काम पूरा है।
\end{proof}

\begin{prop}\ollabel{ordisoidentity}
किन्हीं क्रमसूचकों $\alpha, \beta$ के लिए $\alpha = \beta$ तभी और केवल
तभी जब $\ordeq{\alpha}{\beta}$।
\end{prop}

\begin{proof}
क्रमसूचक सुव्यवस्थाएँ हैं; अतः यह त्रिकोटी
(\olref[basic]{ordtrichotomy}) और \olref[iso]{wellordnotinitial} से तत्काल
मिलता है।
\end{proof}

\begin{prob}\ollabel{probunionordinalsordinal}
सिद्ध कीजिए कि यदि $X$ का प्रत्येक सदस्य क्रमसूचक है, तो $\bigcup X$
क्रमसूचक है।
\end{prob}

\end{document}
